% Source: https://artofproblemsolving.com/community/c3h1881460_inscribed_circle_in_pentagon_geometry_problem
% Stable identity: aops-topic-1881460-q01
Let the pentagon be $ABCDE$ with side $s$, and let $d_1,\dots,d_5$ be the distances from $P$ to the lines $AB,BC,CD,DE,EA$. (For a point near a vertex, the foot of a perpendicular can lie on an extension of a side; the length is still the distance from $P$ to that line.)

\begin{asy}
usepackage("amsmath");
size(6cm);
defaultpen(fontsize(10pt));
pair[] V;for(int k=0;k<5;++k){V.push(dir(90+72*k));}
pair P=(-0.3,-0.2);
pair ft(pair X,pair Y){return X+dot(P-X,Y-X)/dot(Y-X,Y-X)*(Y-X);}
for(int k=0;k<5;++k){draw(P--V[k],gray(0.7)+linewidth(0.4));}
draw(V[0]--V[1]--V[2]--V[3]--V[4]--cycle,linewidth(1));
for(int k=0;k<5;++k){pair X=V[k],Y=V[(k+1)%5];pair F=ft(X,Y);
 real t=dot(F-X,Y-X)/dot(Y-X,Y-X);
 if(t<0) draw(X--F,dotted+linewidth(0.8));
 if(t>1) draw(Y--F,dotted+linewidth(0.8));
 draw(P--F,blue+dashed+linewidth(0.8));
 pair u=unit(Y-X),v=unit(P-F);real s=0.045;
 draw((F+s*u)--(F+s*u+s*v)--(F+s*v),blue+linewidth(0.5));
 dot(F,blue+2.5);}
dot(P);
label("$A$",V[0],N);label("$B$",V[1],W);label("$C$",V[2],SW);label("$D$",V[3],SE);label("$E$",V[4],E);
label("$P$",P+0.13*dir(168));
label("\mbox{}",(min(currentpicture,user=true)+max(currentpicture,user=true))/2,black); // keeps the text after the figure black when LaTeX typesets the labels inline
\end{asy}

(a) Since the pentagon is convex and $P$ is inside it, the segments $PA,\dots,PE$ cut it into the triangles $PAB,PBC,PCD,PDE,PEA$. Each has a base of length $s$ and height equal to the distance from $P$ to that base's line. Adding their areas,
\[
[ABCDE]=\dfrac{1}{2}s\,d_1+\dfrac{1}{2}s\,d_2+\dots+\dfrac{1}{2}s\,d_5=\dfrac {s}{2}\,(d_1+d_2+d_3+d_4+d_5),
\]
so $d_1+\dots+d_5=\dfrac{2[ABCDE]}{s}$, which is the same for every point $P$ inside the pentagon.

(b) Take $P$ to be the centre $O$: then every $d_i$ equals the apothem $a=OM$, where $M$ is the midpoint of $AB$. So the constant is $5a$.

\begin{asy}
usepackage("amsmath");
size(5.5cm);
defaultpen(fontsize(10pt));
pair A=dir(90),B=dir(162),O=(0,0);pair M=(A+B)/2;
draw(dir(90)--dir(162)--dir(234)--dir(306)--dir(18)--cycle,linewidth(1));
draw(O--A,linewidth(0.8));draw(O--B,linewidth(0.8));draw(O--M,dashed+linewidth(0.8));
pair u=unit(A-M),v=unit(O-M);real s=0.06;draw((M+s*u)--(M+s*u+s*v)--(M+s*v));
draw(arc(O,0.42,90,126),linewidth(0.5));
label("$36^\circ$",O+0.6*dir(108));
dot(O);label("$O$",O,S);label("$A$",A,N);label("$B$",B,W);label("$M$",M,NW);
label("$r$",(O+A)/2,E);label("$a$",(O+M)/2+0.09*dir(216));
label("\mbox{}",(min(currentpicture,user=true)+max(currentpicture,user=true))/2,black); // keeps the text after the figure black when LaTeX typesets the labels inline
\end{asy}

In the isosceles triangle $OAB$ the angle at $O$ is $\dfrac{360^\circ}{5}=72^\circ$, and $OM$ bisects it, so $a=r\cos36^\circ$.

To find $\cos36^\circ$, let $\theta=36^\circ$ and $c=\cos\theta$. Since $2\theta=180^\circ-3\theta$, we have $\cos2\theta=-\cos3\theta$, i.e.\ $2c^2-1=-(4c^3-3c)$, so
\[
4c^3+2c^2-3c-1=(c+1)(4c^2-2c-1)=0 .
\]
As $0<c<1$, $c=\dfrac{2+\sqrt{20}}{8}=\dfrac{1+\sqrt5}{4}$. Hence the constant sum is
\[
5r\cos36^\circ=\boxed{\dfrac{5\left(1+\sqrt5\right)r}{4}}.
\]
